% ============================================================
% shared_technical_specification/requirements_table.tex — таблица функциональных
% требований к системе (СКЕЛЕТ; заполняется студентом).
%
% СИСТЕМНАЯ таблица требований командного проекта (отдельный документ,
% не копия личной). Подключается из shared_technical_specification.tex
% (\input{requirements_table}) и shared_test_program_and_methodology.tex
% (\input{../shared_technical_specification/requirements_table}).
%
% Каждое требование объявляется макросом \req{ID}{Текст} в ячейке
% «Требование»: печатается только текст, а панель «Трассировка
% требований» связывает определение со ссылками \reqref{ID} в
% ВКР/ПМИ/РП/РО. Идентификатор в первой ячейке и в \req должен
% совпадать; сохраняйте единый префикс (ТЗ-Ф-01, ТЗ-Ф-02, …).
% ============================================================
\begin{longtable}{|l|>{\raggedright\arraybackslash}p{0.28\textwidth}|>{\raggedright\arraybackslash}p{0.28\textwidth}|>{\centering\arraybackslash}p{0.22\textwidth}|}
\caption{Требования к функциональным характеристикам системы} \label{tab:requirements} \\
\hline
\textbf{№} & \textbf{Пользовательская история} & \textbf{Требование} & \textbf{Область} \\
\hline
\endfirsthead
\multicolumn{4}{l}{Продолжение таблицы \thetable} \\
\hline
\textbf{№} & \textbf{Пользовательская история} & \textbf{Требование} & \textbf{Область} \\
\hline
\endhead
\hline
\endfoot
\hline
\endlastfoot
ТЗ-Ф-01 & Как пользователь, я~хочу \ldots, чтобы \ldots & \req{ТЗ-Ф-01}{Система должна …} & \ldots \\
\hline
ТЗ-Ф-02 & Как пользователь, я~хочу \ldots, чтобы \ldots & \req{ТЗ-Ф-02}{Система должна …} & \ldots \\
\hline
\end{longtable}
